% Chapter 12 of Singular Learning Theory for the Rest of Us: phase transitions of the posterior.
% Local free energies, the posterior's choice between regions, a crossover (the quartic model near
% its singular truth) and a transition (a double-well model).
\newcommand{\phasessig}{\operatorname{sig}}        % the logistic function t -> 1/(1+e^{-t})
\newcommand{\phasesLpop}{\Lambda_{\mathrm{pop}}}   % population crossover curve
\newcommand{\phasesLdat}{\Lambda_{\mathrm{data}}}  % crossover curve of the mean free energy with data
\newcommand{\phasesasinh}{\operatorname{arsinh}}

\sltchapter{Phase Transitions}{How the posterior chooses between singularities}{$y=ax_1+b^2x_2$, $y=abx_1+(2b^4-b^2)x_2$}{ch:phases}
\newtonfoot{0/6,2/2,6/0}{2}{\tfrac12}{a^6+a^2b^2+b^6}
\begin{refsection}[chapters/12-phases.bib]

\begin{chaptersummary}
When several regions of parameter space nearly fit the data, the posterior puts its mass where the local free energy $nL_n(w_\alpha)+\lambda_\alpha\log n$ is smallest, trading fit against learning coefficient. By hand, two models fitted to the same truth $y=\delta x_2$ show two ways the trade plays out: in the quartic model $y=ax_1+b^2x_2$ the effective learning coefficient crosses over from $\frac34$ to $1$ along a universal curve of $n\delta^2$ that dips and overshoots, and in a double-well model the posterior jumps from a segment of misfitting parameters with $\lambda=\frac12$ to two regular true points once $n\delta^2$ reaches about $\log n$. We prove the local free energy formula for models linear in their coefficients, that the posterior chooses the region whose local free energy is smallest by a growing margin and moves toward better fit and larger learning coefficient, and the limiting crossover curve and posterior odds of the two examples. Quadrature over thousands of data sets confirms the curves, including a switch near $n=3000$, where the regular theory, which charges every region the same, gives the better-fitting region most of the posterior already at $n=100$. The theory does not say where the regions come from, the switch point of a single data set is random, and whether training passes through the same regions is open.
\end{chaptersummary}

\section{The question: where does the posterior go?}\label{sec:phases:question}

Train a small network on more and more data and look at where its weights end up. In the toy model of superposition of \textcite{ElhageEtAl2022},\sidenote{\Textcite{ElhageEtAl2022} introduced the model to study how a network can represent more features than it has dimensions, which they call superposition.} a ReLU network asked to store six sparse features in two hidden dimensions, the critical points of the loss include regular polygons: four features stored as a square, five as a pentagon, six as a hexagon; and training curves rest near such configurations before dropping to the next \parencite{ChenEtAl2023}.\index{toy model of superposition} \Textcite{ChenEtAl2023} asked the Bayesian version of the question: as the number $n$ of examples grows, does the posterior move between the same configurations, and at which $n$? A configuration that stores fewer features fits worse, but the weights of the unused feature are partly free, so it is more degenerate. The posterior has to trade fit against degeneracy.

The tool is local. For a measurable set $U\subseteq W$ with positive prior mass, the \term{local free energy}\index{free energy!local} is
\begin{equation}\label{eq:phases:local}
F_n(U)=-\log\int_U\varphi(w)\,e^{-nL_n(w)}\dd w ,
\end{equation}
the free energy of the model restricted to $U$, with the prior not renormalized. The posterior density is $\varphi\,e^{-nL_n}e^{F_n}$, so the posterior mass of $U$ is exactly
\begin{equation}\label{eq:phases:mass}
\Prob(U\mid X^n)=e^{F_n-F_n(U)} .
\end{equation}
The posterior prefers the region with the smallest local free energy. Following physics and \textcite{ChenEtAl2023}, we call a region whose local free energy we follow a \term{phase}\index{phase}, and a change of the region that carries most of the posterior mass a \term{phase transition}\index{phase transition}.\sidenote{In statistical physics a phase transition is a point at which a free energy, in the limit of a large system, fails to be analytic in a parameter such as the temperature. Here the parameter is $n$: the minimum $\min_\alpha(n\Delta_\alpha+\lambda_\alpha\log n)$ of the leading terms of the local free energies has a kink where the minimizing region changes, as a free energy does at a first-order transition. In the statistical mechanics of learning, where the numbers of parameters and of examples grow together, learning curves can jump \parencite[ch.~7, ``Discontinuous learning'']{EngelVandenBroeck2001}; that is not the limit of this chapter.}

\begin{regularanswer}[where the posterior goes]
If the model is regular, then the posterior concentrates in a neighborhood of the maximum-likelihood estimate $\hat w$ of radius of order $n^{-1/2}$ and is asymptotically normal there \parencite[Theorem~10.1]{vanderVaart1998}. Applied region by region, Laplace's approximation charges every region the same $\frac d2\log n$ (\cref{thm:regular:laplace}), so the posterior odds of two regions are their maximized likelihood ratio up to a bounded factor: the region that fits better wins as soon as $n$ times the difference of their losses is large.
\end{regularanswer}

Singular models break the second sentence. Watanabe's formula, applied to the model restricted to a region (\cref{thm:freeenergy:main}), charges it $\lambda_\alpha\log n-(m_\alpha-1)\log\log n$ with its own learning coefficient. A region that fits worse by $\Delta$ per data point but has a smaller learning coefficient keeps the posterior until $n\Delta$ catches up with the difference of the $\lambda_\alpha\log n$. In the double-well model of this chapter at $n=3162$, equal penalties predict that the better-fitting region carries $93\%$ of the posterior; it carries $53\%$. The regular answer also fails more quietly: near a singular truth a region does not have one learning coefficient over the range of $n$ that matters. This chapter asks which region carries the posterior at sample size $n$, when and how sharply it switches, and what replaces the free energy formula when the truth sits close to a singularity.

\section{By hand: one truth, two models}\label{sec:phases:byhand}

Both models of this section have inputs $x=(x_1,x_2)\sim\Normal(0,I_2)$, unit noise, the uniform prior on $[-1,1]^2$, and the truth
\[
y=\delta x_2+\varepsilon,\qquad 0<\delta<1 .
\]
Both are linear in a coefficient vector, $y=g(w)^\top x+\varepsilon$, so $K(w)=\frac12\abs{g(w)-c_0}^2$ with $c_0=(0,\delta)$, as in \cref{ch:freeenergy}. They differ only in how they reach $c_0$.

\paragraph{A crossover.} In the quartic model $y=ax_1+b^2x_2$ the truth is reached at the two points $(0,\pm\sqrt\delta)$; at $\delta=0$ the same model is singular, with $\lambda=\frac34$ (\cref{ch:blowups}). The two answers cannot both hold at every $n$.

\begin{computation}{The quartic model near its singular truth}\label{comp:phases:quartic}
\step The loss $K=\frac12a^2+\frac12(b^2-\delta)^2$ vanishes at $(0,\pm\sqrt\delta)$, where its Hessian is $\diag(1,4\delta)$: both true parameters are regular, and $\lambda=1=\frac d2$.
\step The population free energy $\bar F(n)=-\log\int\varphi e^{-nK}$ factorizes: $\int\varphi e^{-nK}=\frac14I_aI_b$, with $I_a=\int_{-1}^1e^{-na^2/2}\dd a=\sqrt{2\pi/n}\,(1-2\bar\Phi(\sqrt n))$, $\bar\Phi$ the standard normal tail, and $I_b=\int_{-1}^1e^{-n(b^2-\delta)^2/2}\dd b$.
\step Substitute $b=n^{-1/4}u$, the scale of the singular model's posterior. Then $n(b^2-\delta)^2=(u^2-\tau)^2$ with $\tau=\sqrt n\,\delta$, and $I_b=n^{-1/4}\int_{\abs u\le n^{1/4}}e^{-(u^2-\tau)^2/2}\dd u$: the truth enters only through $\tau$.
\step Let $h(s)=\int_\R e^{-(u^2-s)^2/2}\dd u$. Up to the tails beyond $\abs u=n^{1/4}$ and $\abs a=1$, which are exponentially small (\cref{prop:phases:crossover}),
\[
\bar F(n)\approx\tfrac34\log n+\log4-\tfrac12\log(2\pi)-\log h(\sqrt n\,\delta).
\]
\step At $\tau=0$, $h(0)=2^{5/4}\Gamma(\frac54)=2.156$ and $\bar F(n)=\frac34\log n+O(1)$, the singular formula. For large $\tau$ the integrand of $h$ has two Gaussian peaks at $\pm\sqrt\tau$ of width $1/(2\sqrt\tau)$, so $h(\tau)\sim\sqrt{2\pi/\tau}$ and $\bar F(n)=\log n+\frac12\log\delta+\log\frac2\pi+o(1)$, Laplace's formula at the two regular points.
\step The \term{effective learning coefficient}\index{learning coefficient!effective} $\lambda_{\mathrm{eff}}(n)=\dd\bar F/\dd\log n$ is therefore, up to exponentially small terms, a function of $n\delta^2$ alone:
\[
\lambda_{\mathrm{eff}}(n)=\phasesLpop(\sqrt n\,\delta),\qquad\phasesLpop(\tau)=\frac34-\frac\tau2\,\frac{h'(\tau)}{h(\tau)} .
\]
The script \texttt{12-phases-crossover.py} gives $\phasesLpop(0.1)=0.730$, $\phasesLpop(1)=0.803$, $\phasesLpop(10)=1.004$, and $\phasesLpop=\frac78$ at $n\delta^2=1.60$.
\end{computation}

What one had to notice is the scale: the singular posterior of $b$ lives at $n^{-1/4}$, and on that scale the true value $\sqrt\delta$ is invisible until $\sqrt n\,\delta$ is of order one. A coefficient within $O(n^{-1/2})$ of the singular one behaves like the singular point, one much farther away like a regular point, and in between the free energy follows neither formula.\sidenote{The function $h$ is a parabolic cylinder function, $h(s)=\sqrt\pi\,e^{-s^2/4}D_{-1/2}(-s)$, which we checked numerically; nothing below uses the identity.} The passage takes about four decades of $n\delta^2$, and it is not monotone: $\phasesLpop$ dips to $0.708$ at $n\delta^2=0.13$, returns to $\frac34$ where $h$ is largest ($\tau=0.76$), and overshoots to $1.084$ at $n\delta^2=6.7$ (\cref{fig:phases:crossover}).\sidenote{\Textcite{WatanabeAmari2003} studied layered models under this scaling, a truth at Kullback--Leibler distance of order $1/n$ from the singular ones, and found that the Bayes generalization error can then exceed the regular one, an effect of the same kind as the overshoot.} Heuristically, the hump between the wells is responsible: along $a=0$ the loss has a local maximum $\frac12\delta^2$ at $b=0$, so the state density of \cref{ch:volume} is infinite at $\eps=\frac12\delta^2$, and a posterior, which lives where $K$ is of order $1/n$, meets the hump when $\frac12n\delta^2$ is of order one. First the flattened valley lets it spread, then the wells separate and it narrows at the regular rate. This is a \term{crossover}\index{crossover}: one region whose effective learning coefficient changes. The point $(0,0)$ is a saddle of $K$, and nothing separates it from the true points.

\paragraph{A transition.} The \term{double-well model}\index{double-well model}
\[
y=ab\,x_1+q(b)\,x_2+\varepsilon,\qquad q(b)=2b^4-b^2,
\]
reaches $c_0$ differently. Read $b$ as the input weight of a unit, $a$ as its output weight on the first feature, and let the same input weight drive the coefficient $q(b)$ of the second feature: $q$ is $0$ at $b=0$, negative for $0<\abs b<1/\sqrt2$ with minimum $q(\pm\frac12)=-\frac18$, and $q(\pm1)=1$. To fit $\delta x_2$ the unit must pass through weights at which it fits worse than not at all.

\begin{computation}{Two regions of the double-well model}\label{comp:phases:doublewell}
\step The loss is $K=\frac12a^2b^2+\frac12(q(b)-\delta)^2$. It vanishes only where $ab=0$ and $q(b)=\delta$; since $q\le0$ for $\abs b\le1/\sqrt2$, this forces $a=0$ and $b=\pm b_\delta$ with $b_\delta^2=\frac14(1+s)$, $s=\sqrt{1+8\delta}$. There the Jacobian of $g(a,b)=(ab,q(b))$ is $\diag(\pm b_\delta,\pm q'(b_\delta))$, of determinant $D=b_\delta q'(b_\delta)=\frac12s(1+s)$ in absolute value: both true parameters are regular, $\lambda=1$. For $\delta=0.05$, $b_\delta=0.7388$ and $D=1.2916$.
\step Region $A=\{\abs b\le\frac12\}$. There $q(b)\le0$, so $K\ge\frac12\delta^2+\frac12a^2b^2\ge\frac12\delta^2$, with equality exactly on the segment $\{b=0\}$, for every $a$: the unit is off, the misfit is $\Delta=\frac12\delta^2$ per data point, and $a$ is not identified. Across the segment $K-\Delta=(\delta+\frac12a^2)b^2+O(b^4)$, so the segment is a set of minima with nondegenerate curvature across it, and $\lambda=\frac12$, $m=1$ (\cref{prop:blowups:morsebott}).
\step Region $B=\{\abs b\ge\frac12\}$ contains both true parameters. On the ridge $\abs b=\frac12$, $K\ge\frac12(\delta+\frac18)^2$; for $\delta=0.05$ that is $0.0153$ against $\Delta=0.00125$, a barrier of height $\eta=0.0141$.
\step The local population free energies. In $A$, for fixed $a$ the integral over $b$ is $\sqrt{\pi/(n(\delta+a^2/2))}$ to leading order, and $\int_{-1}^1(\delta+\frac12a^2)^{-1/2}\dd a=2\sqrt2\,\phasesasinh(1/\sqrt{2\delta})$, so
\[
\bar F_A(n)=n\Delta+\tfrac12\log n-\log\bigl(\sqrt{\pi/2}\,\phasesasinh(1/\sqrt{2\delta})\bigr)+o(1).
\]
In $B$, Laplace's method gives each true point $\frac14\cdot2\pi/(nD)$, so $\bar F_B(n)=\log n+\log\frac D\pi+o(1)$. For $\delta=0.05$ the constants are $-0.851$ and $-0.889$.
\step The switch: $\bar F_A=\bar F_B$ when $n\Delta=\frac12\log n-0.038$, at $\bar n_c=3198$ ($3232$ from the leading terms alone). With the population loss in place of the data, the posterior mass of $B$ is $0.11$, $0.49$ and $0.94$ at $n=1000$, $3162$ and $5623$.
\end{computation}

What one had to notice is that the misfit turns the segment into a trap: from each of its points, the first move in $b$ makes the coefficient of $x_2$ negative, away from the truth. The segment is the network without the unit, embedded in the larger network as a line of local minima, as \textcite{FukumizuAmari2000} describe for multilayer perceptrons.\sidenote{They show that a critical point of a perceptron with $H-1$ hidden units gives lines of critical points of the perceptron with $H$ units, which can be local minima or saddles, using only the hierarchical structure of the model.} The truth that the quartic model approached through a crossover is now reached by a jump, when the misfit $n\Delta=\frac12\tau^2$ has caught up with the extra $\frac12\log n$; the quartic model needed only $\tau$ of order one. At the same height $\frac12\delta^2$, the singular point of the quartic model is a hilltop and the segment of the double-well model a valley floor (\cref{fig:phases:profiles}). The regular theory charges both regions $\log n$, so it prefers the better-fitting region $B$ at every $n$; with data it gives $B$ most of the posterior already at $n=100$ (\cref{sec:phases:measuring}).

\begin{figure}[t]
\centering
\begin{tikzpicture}
\begin{axis}[splot,height=0.31\linewidth,xmin=-1,xmax=1,ymin=0,ymax=0.03,
  xlabel={$b$},ylabel={$\min_aK(a,b)$},xtick={-1,-0.5,0,0.5,1},
  ytick={0,0.01,0.02,0.03},yticklabel style={/pgf/number format/fixed,/pgf/number format/precision=2},
  legend style={font=\scriptsize,at={(axis cs:0,0.0295)},anchor=north}]
\addplot[theory2,domain=-1:1,samples=401] {0.5*(x^2-0.05)^2};
\addlegendentry{quartic: $\frac12(b^2-\delta)^2$}
\addplot[theory,domain=-1:1,samples=401] {0.5*(2*x^4-x^2-0.05)^2};
\addlegendentry{double well: $\frac12(q(b)-\delta)^2$}
\addplot[guide,domain=-1:1,samples=2,forget plot] {0.00125};
\draw[muted,line width=0.4pt] (axis cs:-0.5,0) -- (axis cs:-0.5,0.03);
\draw[muted,line width=0.4pt] (axis cs:0.5,0) -- (axis cs:0.5,0.03);
\node[font=\scriptsize,anchor=south] at (axis cs:0,0.0075) {region $A$};
\node[font=\scriptsize,anchor=south] at (axis cs:-0.64,0.02) {$B$};
\node[font=\scriptsize,anchor=south] at (axis cs:0.64,0.02) {$B$};
\end{axis}
\end{tikzpicture}
\caption{The loss along the line $a=0$, which is the minimum of $K(a,b)$ over $a$, for the two models at the truth $y=0.05\,x_2$. Dotted: the quartic model, whose singular point $b=0$ is a local maximum at height $\frac12\delta^2$ between the true parameters $\pm0.224$. Solid: the double-well model, whose segment $b=0$ is a valley floor at the same height $\Delta$ (gray), separated from the true parameters $\pm0.739$ by ridges at $b=\pm\frac12$ (vertical lines).}\label{fig:phases:profiles}
\end{figure}

\section{The theorem: local free energies and the posterior's choice}\label{sec:phases:local}

Both computations used the population loss. With data we must show that each local free energy has the expansion the computations suggest, that the posterior then sits in the region with the smallest one, and, in \cref{sec:phases:curves}, the limiting curves of the two examples.

\begin{lemma}[Local free energies]\label{lem:phases:softmin}
Let $U_1,\dots,U_k$ be measurable subsets of $W$ with positive prior mass that cover $W$ and overlap only in sets of prior measure zero. Then
\[
F_n=-\log\sum_{\alpha=1}^ke^{-F_n(U_\alpha)},\qquad \min_\alpha F_n(U_\alpha)-\log k\le F_n\le\min_\alpha F_n(U_\alpha),
\]
and $\Prob(U_\alpha\mid X^n)=e^{-F_n(U_\alpha)}\big/\sum_\beta e^{-F_n(U_\beta)}$ for each $\alpha$.
\end{lemma}

\begin{proof}
The integral of $\varphi e^{-nL_n}$ over $W$ is the sum of its integrals over the $U_\alpha$, because the overlaps carry no prior mass; this is the identity. The sum lies between its largest term $e^{-\min_\alpha F_n(U_\alpha)}$ and $k$ times that term, and the last formula is \eqref{eq:phases:mass}.
\end{proof}

The free energy is a soft minimum of the local ones. We prove the expansion of the local free energies for the models of \cref{thm:freeenergy:linear}, in a region that need not contain the truth.

\begin{assumption}[A region of a linear model]\label{ass:phases:region}
\leavevmode
\begin{enumerate}
\item The data are independent pairs $(x_i,y_i)$. The inputs $x_i\in\R^k$ have the law of $x$, where $\Sigma=\E[xx^\top]$ is positive definite and $\E\abs x^4<\infty$, and the outputs are $y_i=c_0^\top x_i+\varepsilon_i$, with standard normal noise independent of $x_i$.
\item The model is $y=g(w)^\top x+\varepsilon$ with standard normal noise, for a measurable $g\colon W\to\R^k$ on $W\subseteq\R^d$, with a prior density $\varphi$. So $L(w)$ is $\frac12\abs{g(w)-c_0}_\Sigma^2$ plus a constant, where $\abs v_\Sigma^2=v^\top\Sigma v$.
\item The region $U\subseteq W$ is measurable, and $L$ attains its infimum over $U$ at a point $w_U\in U$. Write $K_U=L-L(w_U)\ge0$ on $U$.
\item (Relatively finite variance on $U$.) For some $c_U>0$, $K_U(w)\ge c_U\abs{g(w)-g(w_U)}_\Sigma^2$ for every $w\in U$.
\item (Local volume.) $V_U(\eps)=\int_{U\cap\{K_U<\eps\}}\varphi\,\dd w\sim C_U\eps^{\lambda_U}(\log\frac1\eps)^{m_U-1}$ as $\eps\to0$, for some $C_U>0$, $\lambda_U>0$ and integer $m_U\ge1$.
\end{enumerate}
\end{assumption}

Nothing requires the truth to lie in $U$: the best point $w_U$ may have a misfit $L(w_U)-L(w_0)>0$, as the segment of the double-well model does. Condition 4 is Watanabe's relatively finite variance\index{relatively finite variance} for the model restricted to $U$, with $w_U$ in the role of the truth: for $f_U=\log p(y\mid x,w_U)-\log p(y\mid x,w)$ one computes $\E f_U=K_U(w)$ and $\E f_U^2\ge\abs{g(w)-g(w_U)}_\Sigma^2$, with a matching upper bound when $g$ is bounded. When $U$ contains a true parameter, $K_U=\frac12\abs{g-g(w_U)}_\Sigma^2$ and $c_U=\frac12$. The condition fails when $U$ contains a whole curve of distinct, equally good distributions (\cref{ex:phases:circle}).

\begin{theorem}[Local free energies and the posterior's choice]\label{thm:phases:local}
\begin{enumerate}[label=(\roman*)]
\item Under \cref{ass:phases:region},
\[
F_n(U)=nL_n(w_U)+\lambda_U\log n-(m_U-1)\log\log n+O_p(1).
\]
More precisely, let $R_n=F_n(U)-nL_n(w_U)-\lambda_U\log n+(m_U-1)\log\log n$ and $\kappa_U=-\log(C_U\Gamma(\lambda_U+1))$; let $\Sigma_n=\frac1n\sum_ix_ix_i^\top$, $Z_n=n^{-1/2}\sum_ix_i\varepsilon_i$, $\tilde Z_n=Z_n-\sqrt n(\Sigma_n-\Sigma)(g(w_U)-c_0)$, $\tilde\rho_n^2=\tilde Z_n^\top\Sigma^{-1}\tilde Z_n$; and let $A_n(\eta)$ be the event $(1-\eta)\Sigma\preceq\Sigma_n\preceq(1+\eta)\Sigma$. Then $\Prob(A_n(\eta))\to1$, $\tilde\rho_n^2$ is bounded in probability, and for every $\eta\in(0,\frac12c_U]$ there are numbers $\epsilon_n(\eta)\to0$ such that on $A_n(\eta)$
\begin{equation}\label{eq:phases:band}
\abs{R_n-\kappa_U}\le\frac{2\lambda_U\eta}{c_U}+\frac{\tilde\rho_n^2}{2\eta}+\epsilon_n(\eta).
\end{equation}
\item Let $U_1,\dots,U_k$ be as in \cref{lem:phases:softmin}, each satisfying \cref{ass:phases:region} with its own $w_\alpha$, $\lambda_\alpha$, $m_\alpha$, and put $\Phi_\alpha(n)=nL_n(w_\alpha)+\lambda_\alpha\log n-(m_\alpha-1)\log\log n$. If $\Phi_\beta(n)-\Phi_\alpha(n)\to\infty$ in probability for every $\beta\ne\alpha$, then $\Prob(U_\alpha\mid X^n)\to1$ in probability.
\end{enumerate}
\end{theorem}

\begin{proof}
The plan is the proof of \cref{thm:freeenergy:linear} with the region's best point $w_U$ in the place of the truth: sandwich the random energy on $U$ between two multiples of $nK_U$ and read the two population free energies off the Abelian lemma. Write $\delta_U=g(w_U)-c_0$ and $e(w)=g(w)-g(w_U)$.

\emph{The energy.} Since $y_i-g(w)^\top x_i=\varepsilon_i-(\delta_U+e)^\top x_i$, expanding the squares in $nL_n(w)=\frac12\sum_i(y_i-g(w)^\top x_i)^2+\frac n2\log(2\pi)$ gives $nL_n(w)-nL_n(w_U)=n\,e^\top\Sigma_n\delta_U+\frac n2e^\top\Sigma_ne-\sqrt n\,e^\top Z_n$, while $nK_U(w)=n\,e^\top\Sigma\delta_U+\frac n2e^\top\Sigma e$. Subtracting,
\begin{equation}\label{eq:phases:energy}
nL_n(w)-nL_n(w_U)=nK_U(w)+\tfrac n2\,e^\top(\Sigma_n-\Sigma)e-\sqrt n\,e^\top\tilde Z_n\qquad(w\in U).
\end{equation}

\emph{The noise.} $\tilde Z_n=n^{-1/2}\sum_i\xi_i$ with $\xi_i=x_i\varepsilon_i-(x_ix_i^\top-\Sigma)\delta_U$, independent and identically distributed with mean $0$, because $\varepsilon_i$ is independent of $x_i$ with mean $0$, and finite covariance, because $\E\abs x^4<\infty$. So $\E\tilde\rho_n^2=\E[\xi_1^\top\Sigma^{-1}\xi_1]$ does not depend on $n$, and $\tilde\rho_n^2$ is bounded in probability by Markov's inequality. The law of large numbers gives $\Sigma_n\to\Sigma$ almost surely, hence $\Prob(A_n(\eta))\to1$.

\emph{The sandwich.} On $A_n(\eta)$, $\abs{e^\top(\Sigma_n-\Sigma)e}\le\eta\abs e_\Sigma^2$, and by the Cauchy--Schwarz inequality $\sqrt n\,\abs{e^\top\tilde Z_n}\le\sqrt n\abs e_\Sigma\tilde\rho_n\le\frac\eta2n\abs e_\Sigma^2+\tilde\rho_n^2/(2\eta)$. So the last two terms of \eqref{eq:phases:energy} are at most $\eta n\abs e_\Sigma^2+\tilde\rho_n^2/(2\eta)$ in absolute value, and condition~4 bounds $\eta n\abs e_\Sigma^2$ by $\theta nK_U$, with $\theta=\eta/c_U\le\frac12$. On $A_n(\eta)$, for every $w\in U$,
\[
(1-\theta)\,nK_U(w)-\frac{\tilde\rho_n^2}{2\eta}\le nL_n(w)-nL_n(w_U)\le(1+\theta)\,nK_U(w)+\frac{\tilde\rho_n^2}{2\eta}.
\]

\emph{Integration.} Exponentiating, integrating against $\varphi$ over $U$ and taking $-\log$ gives, with $\bar F_U(s)=-\log\int_U\varphi\,e^{-sK_U}$,
\[
\bar F_U\bigl((1-\theta)n\bigr)-\frac{\tilde\rho_n^2}{2\eta}\le F_n(U)-nL_n(w_U)\le\bar F_U\bigl((1+\theta)n\bigr)+\frac{\tilde\rho_n^2}{2\eta}.
\]

\emph{The Abelian lemma.} \Cref{lem:volume:abelian}, applied to the prior restricted to $U$, whose mass is at most $1$, gives $\int_U\varphi e^{-sK_U}\sim C_U\Gamma(\lambda_U+1)s^{-\lambda_U}(\log s)^{m_U-1}$ as $s\to\infty$ through the real numbers. Since $\log\log(\sigma n)-\log\log n\to0$, $\bar F_U(\sigma n)=\lambda_U\log n-(m_U-1)\log\log n+\lambda_U\log\sigma+\kappa_U+o(1)$ for fixed $\sigma>0$. Take $\sigma=1\pm\theta$, use $\abs{\log(1\pm\theta)}\le2\theta$, and let $\epsilon_n(\eta)$ be the larger $o(1)$ term: this is \eqref{eq:phases:band}. With $\eta=\frac12c_U$ its right side is bounded in probability, and $\Prob(A_n(\frac12c_U))\to1$, so $R_n=O_p(1)$.

\emph{The choice.} By \cref{lem:phases:softmin}, $\Prob(U_\alpha\mid X^n)=\bigl(1+\sum_{\beta\ne\alpha}e^{-(F_n(U_\beta)-F_n(U_\alpha))}\bigr)^{-1}$, and by part~(i) $F_n(U_\beta)-F_n(U_\alpha)=\Phi_\beta(n)-\Phi_\alpha(n)+O_p(1)\to\infty$ in probability.
\end{proof}

When $U$ contains a true parameter, part~(i) is \cref{thm:freeenergy:linear} for the restricted model. When it does not, the leading term is random at the scale $\sqrt n$: by \eqref{eq:phases:energy} with the truth in the place of $w_U$, $nL_n(w)-nS_n=nK(w)+O_p(\sqrt n)$ for each fixed $w$. A misfitting region competes with a true one through fluctuations of that size, which is why the switch point of a single data set is random (\cref{sec:phases:measuring}). For a fixed truth the fluctuations lose in the end.

\begin{corollary}[Large samples]\label{cor:phases:limit}
In the setting of \cref{thm:phases:local}(ii), suppose that some region contains a true parameter, a point $w_0$ with $g(w_0)=c_0$, and put $\Delta_\alpha=L(w_\alpha)-L(w_0)$. Then $\Prob(U_\alpha\mid X^n)\to0$ in probability whenever $\Delta_\alpha>0$; and if, among the regions with $\Delta_\beta=0$, the region $\alpha$ alone has the smallest learning coefficient, or alone the largest multiplicity among those with the smallest learning coefficient, then $\Prob(U_\alpha\mid X^n)\to1$ in probability.
\end{corollary}

\begin{proof}
If $\Delta_\beta=0<\Delta_\alpha$, then $\Phi_\alpha(n)-\Phi_\beta(n)=n\Delta_\alpha+O_p(\sqrt n)+O(\log n)\to\infty$, by the remark above, and $\Prob(U_\alpha\mid X^n)\le e^{-(F_n(U_\alpha)-F_n(U_\beta))}\to0$. If $\Delta_\alpha=\Delta_\beta=0$, then $g(w_\alpha)=g(w_\beta)=c_0$, so $L_n(w_\alpha)=L_n(w_\beta)$ and $\Phi_\beta-\Phi_\alpha=(\lambda_\beta-\lambda_\alpha)\log n-(m_\beta-m_\alpha)\log\log n$, which tends to $\infty$ under the condition on $\alpha$. Part~(ii) finishes the proof.
\end{proof}

Before that limit the posterior moves, and at leading order the order of its moves is forced. Ignoring the fluctuations and the $\log\log n$ terms, each region is a point $(\lambda_\alpha,\Delta_\alpha)$ that pays $n\Delta_\alpha+\lambda_\alpha\log n$.

\begin{proposition}[Fit against complexity]\label{prop:phases:order}
Let $\bar\Phi_\alpha(n)=n\Delta_\alpha+\lambda_\alpha\log n$ for $\alpha=1,\dots,k$, with $\Delta_\alpha\ge0$ and $\lambda_\alpha>0$, and let $e<n_1<n_2$. If $\bar\Phi_\alpha(n_1)<\bar\Phi_\beta(n_1)$ and $\bar\Phi_\beta(n_2)<\bar\Phi_\alpha(n_2)$, then $\Delta_\beta<\Delta_\alpha$ and $\lambda_\beta>\lambda_\alpha$, and on $(e,\infty)$ the two functions cross exactly once, where
\[
\frac n{\log n}=\frac{\lambda_\beta-\lambda_\alpha}{\Delta_\alpha-\Delta_\beta}.
\]
In particular, as $n$ increases through $(e,\infty)$, every change of the unique minimizer lowers $\Delta$ and raises $\lambda$.
\end{proposition}

\begin{proof}
Put $D(n)=(\bar\Phi_\beta(n)-\bar\Phi_\alpha(n))/\log n=(\Delta_\beta-\Delta_\alpha)\,n/\log n+(\lambda_\beta-\lambda_\alpha)$. On $(e,\infty)$ the function $n/\log n$ has derivative $(\log n-1)/\log^2n>0$, so $D$ is strictly monotone unless $\Delta_\beta=\Delta_\alpha$, when it is constant. Since $D(n_1)>0>D(n_2)$, $D$ decreases, so $\Delta_\beta<\Delta_\alpha$, and $\lambda_\beta-\lambda_\alpha=D(n_1)+(\Delta_\alpha-\Delta_\beta)\,n_1/\log n_1>0$; a strictly decreasing function has at most one zero. The last claim is the first applied to consecutive unique minimizers.
\end{proof}

Geometrically, $\bar\Phi_\alpha(n)/n=\Delta_\alpha+t\lambda_\alpha$ with $t=\log n/n$, which decreases from $1/e$ to $0$: the preferred region minimizes a linear function of $(\lambda_\alpha,\Delta_\alpha)$ whose slope shrinks as data arrive. The regions ever preferred lie on the lower boundary of the convex hull of the points and are visited from small $\lambda$ and large misfit toward large $\lambda$ and small misfit; a point above that boundary is never preferred (\cref{ex:phases:hull}).\sidenote{For $c>e$ the equation $n/\log n=c$ has the solution $n=-c\,W_{-1}(-1/c)$ on $(e,\infty)$, with $W_{-1}$ the lower branch of Lambert's function; numerically $n$ lies between $1.27\,c\log c$ and $1.41\,c\log c$ for $c$ between $10^2$ and $10^4$.} \Textcite[App.~C]{ChenEtAl2023} derive the switch equation and show, for transitions that leave the constant terms unchanged, that slopes increase along a sequence of transitions. This is the precise sense in which a growing data set makes the posterior jump to more complex singularities: a region of smaller learning coefficient wins only while its worse fit has not yet been multiplied by enough data. The $\log\log n$ terms, the constants and the fluctuations of $nL_n(w_\alpha)$ shift the switch points, and \cref{sec:phases:curves} computes them in the examples.

\paragraph{The general formula.} For a model that is not linear in its coefficients, the local formula is \cref{thm:freeenergy:main}, in its form for truths that the model does not realize \parencite{Watanabe2010d}, applied to the model restricted to $U$: if the restricted model satisfies its conditions with $p(\cdot\mid w_U)$ as the optimal distribution (among them, $U$ compact and defined by analytic inequalities, $w_U$ interior with the only optimal distribution in $U$, and relatively finite variance with respect to it), then $F_n(U)=nL_n(w_U)+\lambda_U\log n-(m_U-1)\log\log n+O_p(1)$, with $\lambda_U$ and $m_U$ those of $K_U$ on $U$. This is the local free energy formula of \textcite[eq.~(8)]{ChenEtAl2023}, who glue the regions with a partition of unity instead of indicator functions and call the comparison of local free energies \term{internal model selection}\index{internal model selection}.

\begin{observation}[\textcite{ChenEtAl2023}]
In the toy model of superposition with $2$ hidden dimensions and $6$ features, $f(x,w)=\relu(V^\top Vx+c)$ with $V\in\R^{2\times6}$ and $c\in\R^6$ ($d=18$), the losses $\E\norm{x-f(x,w)}^2$ and the learning coefficients of the pentagon ($0.06874$, $\lambda=7$) and the hexagon ($0.04819$, $\lambda=8.5$), with constants from the prior density at the two polygons, predict that the posterior moves from the pentagon to the hexagon at $n\approx600$ (their switch equation gives $601$). Posterior samples from NUTS, with the prior $\Normal(0,I)$, six chains of $5000$ iterations after $500$ of burn-in, and $10$ data sets for each $n$, show the switch between $n=600$ and $n=700$. Below $n\approx400$ the classification of samples becomes unreliable. The authors flag a tension in comparing asymptotic formulas at finite $n$. The same transition was not observed in their stochastic gradient descent runs.
\end{observation}

The study is a preprint, and the agreement is less sharp than it looks. The switch equation charges each polygon $n$ times the loss $\E\norm{x-f}^2$, but the likelihood that their Appendix~F says the sampler used, $p(x\mid w)\propto\exp(-\frac12\norm{x-f(x,w)}^2)$ (their eq.~(3)), has half that energy, an apparent inconsistency by a factor of two. With the halved energy and the same constants the switch is at $n\approx1317$; without the constants it is at $445$, or $1010$ with the halved energy (our computations from their tables); \cref{ch:stages}, which also tries the asymptotic constants, finds predicted switches from about $110$ to about $1300$. Their pentagon with an extra positive bias, which has the hexagon's learning coefficient $8.5$ but the higher loss $0.06180$, is never preferred at leading order.

\section{The theorem: the crossover and the transition}\label{sec:phases:curves}

With data, the computations of \cref{sec:phases:byhand} become limit theorems in which the fluctuations of the data are kept rather than bounded. The tool is a form of dominated convergence for integrals that depend on a few statistics of the data.

\begin{lemma}[Continuous convergence]\label{lem:phases:continuous}
Let $\xi_n\to\xi^*$ in probability for a constant $\xi^*$, and let $\zeta_n$ be bounded in probability. Let $f_n$ and $f$ be measurable, $f$ continuous, with $f_{n_j}(x_j,z_j)\to f(z)$ whenever $n_j\to\infty$, $x_j\to\xi^*$ and $z_j\to z$. Then $f_n(\xi_n,\zeta_n)-f(\zeta_n)\to0$ in probability.
\end{lemma}

\begin{proof}
The plan is to make the convergence uniform on compact sets and then use tightness. Since $\xi_n\to\xi^*$ in probability, there are $r_n\to0$ with $\Prob(\abs{\xi_n-\xi^*}>r_n)\to0$: take $N_j$ increasing with $\Prob(\abs{\xi_n-\xi^*}>1/j)<1/j$ for $n\ge N_j$, and $r_n=1/j$ for $N_j\le n<N_{j+1}$. Fix $M>0$ and let $s_n$ be the supremum of $\abs{f_n(x,z)-f(z)}$ over $\abs{x-\xi^*}\le r_n$, $\abs z\le M$. If $s_n\not\to0$, then there are $\epsilon>0$, $n_j\to\infty$ and such $(x_j,z_j)$ with $\abs{f_{n_j}(x_j,z_j)-f(z_j)}\ge\epsilon$; along a subsequence $z_j\to z$, and then $f_{n_j}(x_j,z_j)\to f(z)$ and $f(z_j)\to f(z)$, a contradiction. Hence
\[
\Prob\bigl(\abs{f_n(\xi_n,\zeta_n)-f(\zeta_n)}>\epsilon\bigr)\le\Prob(\abs{\xi_n-\xi^*}>r_n)+\Prob(\abs{\zeta_n}>M)+\ind\{s_n>\epsilon\},
\]
whose upper limit is at most $\sup_n\Prob(\abs{\zeta_n}>M)$, small for large $M$.
\end{proof}

\begin{proposition}[The crossover]\label{prop:phases:crossover}\index{crossover}
Consider the quartic model $y=ax_1+b^2x_2$ with inputs $\Normal(0,I_2)$, unit noise, the uniform prior on $[-1,1]^2$ and the truth $y=\delta x_2+\varepsilon$, $0\le\delta<1$. Let $h(s)=\int_\R e^{-(u^2-s)^2/2}\dd u$, $\tau=\sqrt n\,\delta$, $Z\sim\Normal(0,1)$, $\Psi(\tau)=-\E\log h(\tau+Z)$, and
\[
\phasesLpop(\tau)=\frac34-\frac\tau2\,\frac{h'(\tau)}{h(\tau)},\qquad\phasesLdat(\tau)=\frac34+\frac\tau2\,\Psi'(\tau).
\]
\begin{enumerate}[label=(\roman*)]
\item (Population.) For every $n\ge1$, $\bar F(n)=\frac34\log n+\log4-\frac12\log(2\pi)-\log h(\tau)+r_n$ with $r_n\ge0$; for each $\bar\delta<1$, $\sup_{\delta\le\bar\delta}r_n$ and $\sup_{\delta\le\bar\delta}\abs{\lambda_{\mathrm{eff}}(n)-\phasesLpop(\tau)}$ tend to $0$ faster than every power of $n$, where $\lambda_{\mathrm{eff}}=\dd\bar F/\dd\log n$.
\item (Data, a truth scaled with $n$.) For fixed $\tau\ge0$ and $\delta=\tau/\sqrt n$, $F_n-nS_n-\frac34\log n$ converges in law to $\log4-\frac12\log(2\pi)-\frac12Z_1^2-\frac12Z_2^2-\log h(\tau+Z_2)$, with $Z_1,Z_2$ independent standard normal, and
\[
\E[F_n-nS_n]=\tfrac34\log n+\log4-\tfrac12\log(2\pi)-1+\Psi(\tau)+o(1).
\]
\item For $r>1$ and fixed $\tau>0$, as $n\to\infty$ with $rn$ an integer and $\delta=\tau/\sqrt n$, the slope $\bigl(\E[F_{rn}-rnS_{rn}]-\E[F_n-nS_n]\bigr)/\log r$ tends to $\frac34+(\Psi(\sqrt r\,\tau)-\Psi(\tau))/\log r$, and this tends to $\phasesLdat(\tau)$ as $r\downarrow1$.
\end{enumerate}
\end{proposition}

\begin{proof}
The plan is to compute (i) exactly, bounding the tails cut off by the box, and to obtain (ii) from the rescaled integral of \cref{comp:phases:quartic} with data, by \cref{lem:phases:continuous}; uniform integrability gives the means, and (iii) follows from (ii).

\emph{(i).} By \cref{comp:phases:quartic}, $r_n=-\log(1-2\bar\Phi(\sqrt n))-\log(1-t_n/h(\tau))\ge0$ with $t_n=\int_{\abs u>n^{1/4}}e^{-(u^2-\tau)^2/2}\dd u$. First, $2\bar\Phi(\sqrt n)\le e^{-n/2}$. For $\abs u\ge n^{1/4}$, $u^2-\tau\ge\sqrt n(1-\delta)>0$, and convexity of the square with $u^2-\sqrt n\ge2n^{1/4}(\abs u-n^{1/4})$ gives $(u^2-\tau)^2\ge n(1-\delta)^2+4n^{3/4}(1-\delta)(\abs u-n^{1/4})$, so $t_n\le e^{-n(1-\delta)^2/2}/(n^{3/4}(1-\delta))$. And $h(\tau)\ge e^{-1/2}/\sqrt{1+\tau}$, since the integrand of $h$ is at least $e^{-1/2}$ on $\sqrt\tau\le\abs u\le\sqrt{\tau+1}$, a set of length at least $1/\sqrt{1+\tau}$; with $\tau\le\sqrt n$ this bounds $r_n$. For the slope, $\lambda_{\mathrm{eff}}(n)=n\E_{\mu_n}[K]$ with $\mu_n\propto e^{-nK}\varphi$, a product of laws of $a$ and $b$. The law of $a$ contributes $\frac12\E[v^2\mid\abs v\le\sqrt n]$ for standard normal $v$, within $\sqrt n\,e^{-n/2}$ of $\frac12$ for $n\ge2$. With $u=n^{1/4}b$ and $f(u)=\frac12(u^2-\tau)^2$, the law of $b$ contributes the mean of $f$ under the density proportional to $e^{-f}$, conditioned on $\abs u\le n^{1/4}$. Without the conditioning, integration by parts gives $\E[uf'(u)]=1$, and $uf'(u)=4f(u)+2\tau(u^2-\tau)$ with $\E[u^2-\tau]=h'(\tau)/h(\tau)$, so $\frac12+\E f=\phasesLpop(\tau)$. The conditioning moves $\E f$ by at most $(\E[f;\abs u>n^{1/4}]+\E[f]\Prob(\abs u>n^{1/4}))/\Prob(\abs u\le n^{1/4})$, which the same tail bound, applied to $f/2$ through $fe^{-f}\le e^{-f/2}$, makes smaller than every power of $n$.

\emph{(ii), the limit law.} By \eqref{eq:phases:energy} with the truth in the place of $w_U$, $nK_n(w)=\frac n2e^\top\Sigma_ne-\sqrt n\,e^\top Z_n$ with $e=(a,b^2-\delta)$. Substituting $a=v/\sqrt n$ and $b=n^{-1/4}u$ makes $\sqrt n\,e=\zeta=(v,u^2-\tau)$. With $D_n=[-\sqrt n,\sqrt n]\times[-n^{1/4},n^{1/4}]$,
\[
F_n-nS_n-\tfrac34\log n=\log4-\log\int_{D_n}e^{-\frac12\zeta^\top\Sigma_n\zeta+\zeta^\top Z_n}\dd v\,\dd u=:f_n(\Sigma_n,Z_n).
\]
Let $f(z)$ be the same expression with $\Sigma_n=I$ and $D_n=\R^2$; completing the squares in $v$ and in $u^2-\tau$ gives $f(z)=\log4-\frac12\log(2\pi)-\frac12z_1^2-\frac12z_2^2-\log h(\tau+z_2)$, which is continuous. If $S_j\to I$, $z_j\to z$ and $n_j\to\infty$, then for large $j$ the integrand is at most $\exp(-\frac14\abs\zeta^2+(\abs z+1)\abs\zeta)\le\exp(-\frac18\abs\zeta^2+2(\abs z+1)^2)$, integrable over $\R^2$, and dominated convergence gives $f_{n_j}(S_j,z_j)\to f(z)$. Since $\Sigma_n\to I$ in probability and $Z_n$ has mean $0$ and covariance $I$, \cref{lem:phases:continuous} gives $F_n-nS_n-\frac34\log n-f(Z_n)\to0$ in probability; the central limit theorem, continuous mapping and Slutsky's lemma (\cref{def:toolkit:convergence}) give the limit law.

\emph{(ii), the means.} With $\rho_n^2=Z_n^\top\Sigma_n^{-1}Z_n$, Cauchy--Schwarz bounds the exponent above by $-\frac14\lambda_{\min}(\Sigma_n)\abs\zeta^2+\rho_n^2$ and below by $-\lambda_{\max}(\Sigma_n)\abs\zeta^2-\frac12\rho_n^2$. Integrating the upper bound over $\R^2$, using $\int e^{-s\abs\zeta^2}=\sqrt{\pi/s}\,(2s)^{-1/4}h(\sqrt{2s}\,\tau)$ and that $h$ is bounded, and the lower bound over $[-1,1]^2\subseteq D_n$,
\[
\bigl|F_n-nS_n-\tfrac34\log n\bigr|\le C\bigl(1+\rho_n^2+\lambda_{\max}(\Sigma_n)+\abs{\log\lambda_{\min}(\Sigma_n)}\bigr).
\]
Each term has a second moment bounded in $n$: $\rho_n^2$ is $\chi^2_2$ given the inputs; $\lambda_{\max}(\Sigma_n)\le\tr\Sigma_n$; and $\lambda_{\min}(\Sigma_n)\ge\det\Sigma_n/\tr\Sigma_n$, where $n^2\det\Sigma_n$ is a product of independent $\chi^2_n$ and $\chi^2_{n-1}$ variables by Bartlett's decomposition. So the sequence is uniformly integrable, and its mean converges to $\E f(Z)$.

\emph{(iii).} At sample size $rn$ the same truth is $\delta=\sqrt r\,\tau/\sqrt{rn}$, so (ii) applies with $\sqrt r\,\tau$; subtract. Since $h'/h$ grows at most linearly, $\Psi'(\tau)=-\E[(h'/h)(\tau+Z)]$ by dominated convergence, and the difference quotient tends to $\frac34+\frac\tau2\Psi'(\tau)$.
\end{proof}

With data the free energy depends on $\tau+Z_2$: the noise in the coefficient of $x_2$ has the same order $n^{-1/2}$ as the truth and blurs the hump, so $\phasesLdat$ dips lower and overshoots less than $\phasesLpop$, with a minimum $0.675$ at $n\delta^2=0.29$, a maximum $1.067$ at $n\delta^2=12.8$, and the value $\frac78$ at $n\delta^2=3.1$.

\begin{proposition}[The transition]\label{prop:phases:doublewell}\index{phase transition}
Consider the double-well model $y=abx_1+q(b)x_2$, $q(b)=2b^4-b^2$, with inputs $\Normal(0,I_2)$, unit noise, the uniform prior on $[-1,1]^2$ and the truth $y=\delta x_2+\varepsilon$, $0<\delta<1$. Let $U_A=\{\abs b\le\frac12\}$, $U_B=\{\abs b\ge\frac12\}$, $S_{xx}=\sum_ix_ix_i^\top$, $S_{xy}=\sum_ix_iy_i$, $\kappa_B=\log(D/\pi)$ with $D$ as in \cref{comp:phases:doublewell}, and
\[
\kappa_A(z)=-\log\Bigl(\frac14\int_{-1}^1\sqrt{\frac\pi{\delta+a^2/2}}\,\exp\Bigl(\frac{a^2z^2}{4\delta+2a^2}\Bigr)\dd a\Bigr).
\]
\begin{enumerate}[label=(\roman*)]
\item $F_n(U_A)=nL_n(0)+\frac12\log n+\kappa_A\bigl((S_{xy})_1/\sqrt n\bigr)+o_p(1)$, with $L_n(0)$ the loss on the segment $b=0$.
\item $F_n(U_B)=nL_n(0)-\frac12S_{xy}^\top S_{xx}^{-1}S_{xy}+\log n+\kappa_B+o_p(1)$.
\item With $\omega_n=\frac12S_{xy}^\top S_{xx}^{-1}S_{xy}-\frac12\log n+\kappa_A\bigl((S_{xy})_1/\sqrt n\bigr)-\kappa_B$,
\[
\log\frac{\Prob(U_B\mid X^n)}{\Prob(U_A\mid X^n)}=\omega_n+o_p(1),
\]
so $\Prob(U_B\mid X^n)-\phasessig(\omega_n)\to0$ in probability, where $\phasessig(t)=1/(1+e^{-t})$.
\end{enumerate}
\end{proposition}

\begin{proof}
The plan is to obtain (ii) from Laplace's approximation for regular models, (i) from \cref{lem:phases:continuous} with an explicit dominating bound, and (iii) as their difference.

\emph{Region $B$.} The map $(a,b)\mapsto(-a,-b)$ preserves the prior and $g$, so $F_n(U_B)=F_n(U_B^+)-\log2$ with $U_B^+=[-1,1]\times[\frac12,1]$. On $U_B^+$, conditions (R1)--(R6) of \cref{thm:regular:laplace} hold for the prior renormalized to $U_B^+$, which is uniform with density $1$: $L$ has the unique minimizer $(0,b_\delta)$, interior because $\frac12<b_\delta<1$, with Hessian $Dg^\top Dg$ of determinant $D^2$, and the log-likelihood is a polynomial in $w$ with coefficients of finite moments. Since $U_B^+$ has prior mass $\frac14$, $F_n(U_B^+)=nL_n(\hat w)+\log\frac n{2\pi}+\log D+\log4+o_p(1)$, with $\hat w$ a minimizer of $L_n$ over $U_B^+$. As $g$ maps a neighborhood of $(0,b_\delta)$ diffeomorphically onto a neighborhood of $c_0$ and $S_{xx}^{-1}S_{xy}\to c_0$ in probability, with probability tending to one $nL_n(\hat w)$ is the minimum over all $c\in\R^2$ of $nL_n(0)+\frac12c^\top S_{xx}c-c^\top S_{xy}$, which is $nL_n(0)-\frac12S_{xy}^\top S_{xx}^{-1}S_{xy}$. This gives (ii).

\emph{Region $A$.} On $U_A$, $g=(ab,q(b))$ vanishes on the segment, so $nL_n(w)-nL_n(0)=\frac12g^\top S_{xx}g-g^\top S_{xy}$. Substitute $b=u/\sqrt n$ and write $q_n(u)=q(u/\sqrt n)=-u^2/n+2u^4/n^2$; then
\[
F_n(U_A)-nL_n(0)-\tfrac12\log n=\log4-\log\int_{-1}^1\int_{\abs u\le\sqrt n/2}e^{-H_n(a,u)}\dd u\,\dd a,
\]
\begin{multline*}
H_n=\frac{(S_{xx})_{11}}{2n}\,a^2u^2+\frac{(S_{xx})_{12}}{\sqrt n}\,au\,q_n(u)+\frac{(S_{xx})_{22}}{2n}\,n\,q_n(u)^2\\
-\frac{(S_{xy})_1}{\sqrt n}\,au-\frac{(S_{xy})_2}n\,n\,q_n(u).
\end{multline*}
Let $\xi_n=(\Sigma_n,(S_{xy})_2/n)$, which tends to $(I,\delta)$ in probability, and $\zeta_n=(\sqrt n(\Sigma_n)_{12},\tilde Z_n)$ with $\tilde Z_n=(S_{xy}-nc_0)/\sqrt n$, which is bounded in probability by the central limit theorem; since $c_0=(0,\delta)$, the first coordinate of $\tilde Z_n$ is $(S_{xy})_1/\sqrt n$. Because $q_n(u)\to0$ and $nq_n(u)\to-u^2$, $H_n(a,u)\to(\delta+\frac12a^2)u^2-zau$ along statistics with $\xi\to(I,\delta)$ and $(S_{xy})_1/\sqrt n\to z$. For the domination, let the coefficients $\sqrt n(\Sigma_n)_{12}$ and $(S_{xy})_1/\sqrt n$ of $H_n$ be at most $T$ in absolute value, $(S_{xy})_2/n\ge\frac12\delta$, and $(\Sigma_n)_{11},(\Sigma_n)_{22}\ge0$, as near the limits. On $\abs u\le\frac12\sqrt n$, $-nq_n(u)=u^2(1-2u^2/n)\ge\frac12u^2$ and $\abs{u\,q_n(u)}\le\abs u^3/n\le u^2/(2\sqrt n)$; the first and third terms of $H_n$ are nonnegative, so $H_n\ge\frac14\delta u^2-Tu^2/(2\sqrt n)-T\abs u\ge\frac18\delta u^2-T\abs u$ once $\sqrt n\ge4T/\delta$, a bound integrable over $[-1,1]\times\R$. Dominated convergence verifies the hypothesis of \cref{lem:phases:continuous}, with the continuous limit $\log4-\log\int_{-1}^1\int_\R e^{-(\delta+a^2/2)u^2+zau}\dd u\,\dd a=\kappa_A(z)$, and the lemma gives (i).

\emph{The odds.} By \cref{lem:phases:softmin} the log odds are $F_n(U_A)-F_n(U_B)$; subtract (ii) from (i), and use that $\phasessig$ is Lipschitz.
\end{proof}

Here $\frac12S_{xy}^\top S_{xx}^{-1}S_{xy}=nL_n(0)-\min_cnL_n(c)$ is the log-likelihood ratio of the least-squares fit against the zero function, and $\omega_n$ is that ratio minus the singular penalty $\frac12\log n$, plus constants; the regular theory charges no penalty. With the population statistics $S_{xx}=nI$, $S_{xy}=nc_0$, $\omega_n=n\Delta-\frac12\log n+\kappa_A(0)-\kappa_B$, whose zero is $\bar n_c$. With data the ratio fluctuates around $n\Delta$ with standard deviation about $\sqrt{2n\Delta}$, which is $2.7$ at $n=3000$ and $\delta=0.05$, comparable to the penalty $\frac12\log n=4$.

\section{Measuring it}\label{sec:phases:measuring}

\paragraph{The crossover.} The script \texttt{12-phases-crossover.py} measures the effective learning coefficient of the quartic model with data, for $\delta=10^{-2}$, $10^{-3}$, $10^{-4}$. For each $\delta$ it draws $3000$ data sets through their sufficient statistics, $S_{xx}$ from the Wishart law by Bartlett's decomposition\sidenote{A $2\times2$ Wishart matrix with $n$ degrees of freedom is $LL^\top$ with $L$ lower triangular and independent entries, $L_{11}^2\sim\chi^2_n$, $L_{22}^2\sim\chi^2_{n-1}$ and $L_{21}\sim\Normal(0,1)$; nested data sets add independent Wishart increments.} and $\sum_ix_i\varepsilon_i\sim\Normal(0,S_{xx})$ given it, and grows each data set by doubling $n$, from the larger of $100$ and $e^{-6}/\delta^2$ to $n\delta^2\approx e^{6.5}$ ($17$ to $19$ sizes, up to $n=6.5\times10^{10}$). Each free energy is a one-dimensional integral: for fixed $b$ the energy is quadratic in $a$, so the integral over $a\in[-1,1]$ is a Gaussian integral with error functions, and the integral over $b$ a trapezoid rule on $1200$ points of the scale $u=n^{1/4}b$, which doubling changes by less than $2\times10^{-13}$. The measured coefficient between consecutive sizes is the difference of the means of $F_n-nS_n+\frac12\rho_n^2-1$ divided by $\log2$; since $\rho_n^2$ is $\chi^2_2$ given the inputs, this control variate leaves the mean unchanged and cuts the standard deviation per data set from $1.33$ to about $0.45$. The intervals are Student-$t$ $95\%$ intervals over the paired differences (\cref{sec:toolkit:intervals}). The run takes $50$ to $75$ seconds.

\begin{figure}[t]
\centering
\begin{tikzpicture}
\begin{axis}[splot,height=0.44\linewidth,xmin=-6.5,xmax=7,ymin=0.6,ymax=1.13,
  xlabel={$\log(n\delta^2)$},ylabel={effective learning coefficient},
  xtick={-6,-4,-2,0,2,4,6},ytick={0.6,0.7,0.75,0.8,0.9,1,1.1},
  legend pos=north west,legend style={font=\scriptsize}]
\addplot[regularpred,domain=-6.5:7,samples=2] {1};
\addlegendentry{regular: $\frac d2=1$}
\addplot[theory] table[x=x,y=data] {figures/data/12-phases-crossover-theory.dat};
\addlegendentry{singular, with data: $\phasesLdat$}
\addplot[theory2] table[x=x,y=pop] {figures/data/12-phases-crossover-theory.dat};
\addlegendentry{singular, population: $\phasesLpop$}
\addplot[guide,domain=-6.5:7,samples=2,forget plot] {0.75};
\addplot[measured,ebars] table[x=x,y=lam,y error plus=hi,y error minus=lo] {figures/data/12-phases-crossover-d2.dat};
\addlegendentry{$\delta=10^{-2}$}
\addplot[measured2,ebars] table[x=x,y=lam,y error plus=hi,y error minus=lo] {figures/data/12-phases-crossover-d3.dat};
\addlegendentry{$\delta=10^{-3}$}
\addplot[measured3,ebars] table[x=x,y=lam,y error plus=hi,y error minus=lo] {figures/data/12-phases-crossover-d4.dat};
\addlegendentry{$\delta=10^{-4}$}
\end{axis}
\end{tikzpicture}
\caption{The effective learning coefficient of the quartic model at the truth $y=\delta x_2$ against $\log(n\delta^2)$; inputs $\Normal(0,I_2)$, unit noise, uniform prior on $[-1,1]^2$. Points: slopes of the mean normalized free energy between $n$ and $2n$, over $3000$ nested data sets for each $\delta$ (seed $1200$), by exact quadrature, with Student-$t$ $95\%$ intervals. Solid bronze: $\phasesLdat$ of \cref{prop:phases:crossover}. Dotted bronze: $\phasesLpop$, which the population free energies of the three values of $\delta$ reproduce to $5\times10^{-8}$. Dashed slate: the regular value $1$. Gray: $\frac34$, the learning coefficient at $\delta=0$. Script: \texttt{12-phases-crossover.py}.}\label{fig:phases:crossover}
\end{figure}

In \cref{fig:phases:crossover} the three series fall on one curve: over eight orders of magnitude of $n$, the effective learning coefficient depends on $n$ and $\delta$ only through $n\delta^2$, and it follows the curve with data, not the population curve. Of the $52$ measured slopes, $51$ lie within their intervals of the prediction of \cref{prop:phases:crossover}(iii) for doubling, which differs from $\phasesLdat$ by at most $0.004$; the exception, $0.656\pm0.023$ against $0.680$ at $n\delta^2=0.45$, is about what $52$ intervals at $95\%$ lead one to expect. The dip and the overshoot are measured: the smallest slopes of the three series are $0.656$, $0.666$, $0.672$ and the largest $1.064$, $1.065$, $1.067$, against the predicted extremes $0.675$ and $1.067$. For $n\delta^2<0.01$ the coefficient is within $0.03$ of $\frac34$, for $n\delta^2>100$ within $0.01$ of $1$; in between, which for $\delta=10^{-3}$ means $10^4\lesssim n\lesssim10^8$, neither formula describes the free energy.

\paragraph{The transition.} The script \texttt{12-phases-transition.py} measures the posterior mass of region $B$ of the double-well model with $\delta=0.05$, for $500$ nested data sets (seed $1210$) at $25$ sample sizes from $10^2$ to $10^5$, equally spaced in $\log n$. It integrates over $a$ in closed form and over $b$ by the trapezoid rule on $8001$ points of $[-1,1]$, the boundaries $\pm\frac12$ among them (doubling changes the logarithms of the region integrals by less than $6\times10^{-8}$). The prediction of \cref{prop:phases:doublewell} is the mean of $\phasessig(\omega_n)$ over $10^5$ independent draws of $(S_{xx},S_{xy})$ for each $n$; the regular prediction is the mean of $\phasessig$ of the maximized log-likelihood ratio of the two regions, the odds that Laplace's approximation gives when both regions pay $\log n$. The run takes $55$ to $95$ seconds.

\begin{figure}[t]
\centering
\begin{tikzpicture}
\begin{axis}[splot,height=0.44\linewidth,xmode=log,xmin=80,xmax=1.25e5,ymin=-0.02,ymax=1.02,
  xlabel={sample size $n$},ylabel={posterior mass of region $B$},
  legend pos=south east,legend style={font=\scriptsize}]
\addplot[regularpred] table[x=n,y=reg] {figures/data/12-phases-transition.dat};
\addlegendentry{regular: equal penalties}
\addplot[theory] table[x=n,y=pred] {figures/data/12-phases-transition.dat};
\addlegendentry{singular: $\E[\phasessig(\omega_n)]$}
\addplot[theory2] table[x=n,y=pop] {figures/data/12-phases-transition.dat};
\addlegendentry{population loss}
\addplot[measured,ebars] table[x=n,y=pB,y error plus=hi,y error minus=lo] {figures/data/12-phases-transition.dat};
\addlegendentry{measured}
\addplot[guide,forget plot] coordinates {(3198,-0.02) (3198,1.02)};
\end{axis}
\end{tikzpicture}
\caption{The posterior mass of the regular region $B=\{\abs b\ge\frac12\}$ of the double-well model at the truth $y=0.05\,x_2$, against $n$; inputs $\Normal(0,I_2)$, unit noise, uniform prior on $[-1,1]^2$. Dots: means over $500$ nested data sets (seed $1210$) of the exact posterior mass, by quadrature, with Student-$t$ $95\%$ intervals. Solid bronze: the prediction $\E[\phasessig(\omega_n)]$ of \cref{prop:phases:doublewell}. Dotted bronze: the posterior mass with the population loss, which switches at $\bar n_c=3198$ (gray line). Dashed slate: the regular prediction, the mean of $\phasessig$ of the maximized log-likelihood ratio of the regions. Script: \texttt{12-phases-transition.py}.}\label{fig:phases:transition}
\end{figure}

In \cref{fig:phases:transition}, from $n\approx400$ on, the measured masses and the prediction agree within the intervals at all $20$ sample sizes; at $n=3162$ the measured mass is $0.534\pm0.031$ and the prediction $0.542$. Below $n\approx400$ the prediction is low by $0.03$ to $0.045$, the $o_p(1)$ of \cref{prop:phases:doublewell} made visible: the log odds of the quadrature exceed $\omega_n$ by $0.25$ on average at $n=100$, $0.044$ at $n=1000$ and $0.005$ at $n=10^4$ (standard deviations $0.33$, $0.084$, $0.026$). The measured curve first falls, from $0.26$ at $n=100$ to $0.20$ at $n=422$, because at $n=0$ the posterior is the prior, which gives each region half its mass: the first move of the posterior is out of the prior and into the singular region. It then rises through $\frac12$ near $n=3000$ and reaches $0.97$ at $n=10^4$. The regular prediction is $0.62$ already at $n=100$ and $0.93$ at $n=3162$.

Each data set switches sharply, but at its own $n$. The fraction of data sets that give $B$ more than half the posterior mass is $0.17$, $0.51$ and $0.83$ at $n=1000$, $3162$ and $5623$, against $0.18$, $0.53$ and $0.81$ for the fraction of draws with $\omega_n>0$, whose median is at $n=2972$; the averaged transition is therefore much wider than the population one, which goes from $0.19$ to $0.94$ between $n=1778$ and $5623$. The ridge separates the regions once $n\eta$ is large: the posterior mass of the band $0.4\le\abs b\le0.6$, averaged over the data sets, is $8.8\%$ at $n=100$, $1.3\times10^{-4}$ at $n=1000$ and $3\times10^{-12}$ at $n=3162$.

\begin{figure}[t]
\centering
\begin{tikzpicture}
\begin{groupplot}[group style={group size=3 by 1,horizontal sep=4mm,yticklabels at=edge left},
  splot,width=0.37\linewidth,height=0.35\linewidth,xmin=-1,xmax=1,ymin=-1,ymax=1,
  xtick={-1,0,1},ytick={-1,-0.5,0,0.5,1},xlabel={$a$},title style={yshift=-1ex}]
\nextgroupplot[ylabel={$b$},title={$n=300$}]
\addplot[guide,domain=-1:1,samples=2] {0.5};
\addplot[guide,domain=-1:1,samples=2] {-0.5};
\addplot[draws] table[x=a,y=b] {figures/data/12-phases-draws-n1.dat};
\nextgroupplot[title={$n=3000$}]
\addplot[guide,domain=-1:1,samples=2] {0.5};
\addplot[guide,domain=-1:1,samples=2] {-0.5};
\addplot[draws] table[x=a,y=b] {figures/data/12-phases-draws-n2.dat};
\nextgroupplot[title={$n=30000$}]
\addplot[guide,domain=-1:1,samples=2] {0.5};
\addplot[guide,domain=-1:1,samples=2] {-0.5};
\addplot[draws] table[x=a,y=b] {figures/data/12-phases-draws-n3.dat};
\end{groupplot}
\end{tikzpicture}
\caption{Draws from the posterior of the double-well model for one data set at the truth $y=0.05\,x_2$, at $n=300$, $3000$ and $30000$ (nested; seed $1221$), $2000$ draws each: $b$ from its marginal by the inverse distribution function on $40001$ points, then $a$ from its conditional law, a truncated normal. Gray: the ridges $b=\pm\frac12$. The posterior mass of region $B$ is $0.15$, $0.49$ and $1.00$. Script: \texttt{12-phases-transition.py}.}\label{fig:phases:draws}
\end{figure}

\Cref{fig:phases:draws} shows one data set before, at and after its switch. At $n=300$ most draws lie along the segment $b=0$, over all values of $a$, thicker where $a$ is small because the valley across the segment has curvature $2\delta+a^2$. At $n=3000$ half the mass has moved to two small clusters at $(0,\pm b_\delta)$, with no draws in between; at $n=30000$ the segment is empty. The draws in $B$ all have $a$ near $0$: the true parameters lie on the other arm, $a=0$, of the product model's cross.

\section{The ledger}\label{sec:phases:ledger}

Where the posterior goes is decided by local free energies, which for models linear in their coefficients are $nL_n(w_\alpha)+\lambda_\alpha\log n-(m_\alpha-1)\log\log n$ up to a bounded term, whether or not the region contains the truth. At leading order a growing sample moves the posterior along the lower convex hull of the points $(\lambda_\alpha,\Delta_\alpha)$, toward better fit and larger learning coefficient. Near a singular truth the effective learning coefficient is a universal, non-monotone function of $n\delta^2$; between separated regions the posterior odds are a likelihood ratio minus the difference of the $\lambda\log n$ terms, with a fluctuation of order $\sqrt n$ that makes each data set switch at its own $n$.

\ledgercard{the posterior mass of a region $U$, through its local free energy $F_n(U)$}{every region pays $\frac d2\log n$; the posterior sits in the region of the maximum-likelihood estimate once $n$ times the loss gap is large}{$F_n(U_\alpha)=nL_n(w_\alpha)+\lambda_\alpha\log n-(m_\alpha-1)\log\log n+O_p(1)$, and the posterior takes the smallest (\cref{thm:phases:local}); transitions in $n$ lower the loss and raise $\lambda$ (\cref{prop:phases:order})}{the quartic model at $y=\delta x_2$: $\lambda_{\mathrm{eff}}$ from $\frac34$ to $1$ along a universal curve of $n\delta^2$, dipping to $0.68$ and overshooting to $1.07$; the double well at $\delta=0.05$: from a segment with $\lambda=\frac12$ to regular points with $\lambda=1$ at $n\approx3000$}{$51$ of $52$ slopes within their intervals; posterior mass of $B$ $0.24$, $0.53$, $0.97$ at $n=10^3$, $3162$, $10^4$, against $0.25$, $0.54$, $0.96$ predicted and $0.75$, $0.93$, $1.00$ from the regular theory}{a fixed model and truth, or a truth scaled with $n$; regions with relatively finite variance and known local volumes; linear models in the proofs, Watanabe's conditions in general; exact posteriors; transitions in $n$, not in training time}

\section{Where the theory stops}\label{sec:phases:stops}

\paragraph{What was assumed.} The local formula is asymptotic for fixed regions and a fixed truth, and not uniform in the truth: in the quartic model it holds with $\lambda=1$ for every $\delta>0$, yet for $n\delta^2\lesssim100$ the free energy grows at a rate between $0.68$ and $1.07$ instead, a mispricing over a range of $n$ that grows like $\delta^{-2}$. The regions are an input: the theorem does not say how to find them in a network, and \textcite{ChenEtAl2023} define theirs by classifying posterior samples. The proofs cover models linear in their coefficients; in general each region needs Watanabe's conditions, including analyticity, which the ReLU of the toy model of superposition violates. Everything concerns the exact posterior at inverse temperature $1$, as a function of $n$; nothing here is about training time.

\paragraph{Where it breaks.} The terms below $n\Delta+\lambda\log n$ are not small where transitions happen. In the double-well model the constants nearly cancel ($3198$ against $3232$), but in the toy model of superposition the constants and the convention for the energy move the predicted switch from about $110$ to about $1300$ (\cref{ch:stages}); and in \cref{ch:coefficient} a $\log\log n$ term keeps a mixed prior's posterior in the region of larger learning coefficient until $n\approx10^{13}$. The fluctuation of order $\sqrt n$ makes the switch of a single data set random: in \cref{fig:phases:transition} the fraction of data sets that have switched rises from $0.17$ to $0.83$ while $n$ grows by a factor of $5.6$. In a preprint on transformers trained for in-context linear regression, \textcite{CarrollEtAl2025} read the move from a ridge-regression solution to one fitted to the finite set of training tasks as a trade-off of this kind; above a threshold of task diversity the move does not happen, although the Bayesian account predicts it, and they leave the reason open, from capacity and under-training to a real difference between Bayesian and dynamical selection. Finally, the posterior mass of a region is easy to define and hard to sample: a ridge of height $\eta$ holds a local sampler for a time that grows like $e^{n\eta}$, and at the switch of the double-well model $n\eta\approx45$ (\cref{ex:phases:mcmc}); our figures use quadrature.\sidenote{Exchange Monte Carlo, a method from statistical physics that swaps chains run at several temperatures, was brought to singular learning machines by \textcite{NagataWatanabe2008a} because Metropolis sampling of their posteriors is costly.}\index{exchange Monte Carlo}

\paragraph{Open problems.} \Textcite{ChenEtAl2023} conjecture that the transitions met during training have Bayesian antecedents, transitions of the posterior in $n$ between the same regions;\index{Bayesian antecedent hypothesis} what links the two, and when training follows the order of \cref{prop:phases:order}, is open as of 2026, and \cref{ch:stages} returns to it. The dynamics of gradient learning near singularities, which produce plateaus \parencite{FukumizuAmari2000,AmariEtAl2006,WeiEtAl2008,CousseauEtAl2008}, are understood for small networks; their relation to the posterior's phases is not. Uniform asymptotics near singular truths are open beyond examples and the layered models of \textcite{WatanabeAmari2003}: which scaling variable replaces $n\delta^2$ for a family of truths approaching a singular one, and is the crossover function fixed by the geometry of the singularity? The lower-order terms of local free energies, which decide transitions at practical $n$, are unknown in general, and \textcite{CarrollEtAl2025} list them among the open questions of the loss--complexity account.

\begin{exercises}
\item \consolidation (The crossover curve.) Prove that $\E_\tau[\frac12(u^2-\tau)^2]=\frac14-\frac\tau2h'(\tau)/h(\tau)$ for the density proportional to $e^{-(u^2-\tau)^2/2}$; show that $\phasesLpop(0)=\frac34$ by identifying the law of $\frac12u^4$ at $\tau=0$ as a Gamma law with shape $\frac14$, and that for $\tau>0$, $\phasesLpop(\tau)<\frac34$ exactly where $h'(\tau)>0$. Check numerically that $h$ is largest at $\tau=0.765$.\label{ex:phases:curve}

\item \consolidation (The other side of the singularity.) Fit the quartic model to $y=-\delta x_2+\varepsilon$, which it cannot realize. Show that the best parameter is $(0,0)$, with misfit $\frac12\delta^2$, and that $\bar F(n)-\frac12n\delta^2$ has effective coefficient $\frac34+\frac\tau2\E[u^2]$ under the density proportional to $e^{-\tau u^2-u^4/2}$, which tends to $1$; check numerically that, unlike $\phasesLpop$, it increases monotonically ($0.773$, $0.895$, $0.996$ at $\tau=0.1$, $1$, $10$).\label{ex:phases:otherside}

\item \consolidation (Lower convex hull.) Show that a region that is the unique preferred one at leading order for some $n>e$ has its point $(\lambda_\alpha,\Delta_\alpha)$ at a vertex of the lower boundary of the convex hull of all the points, and deduce the statement of \textcite[App.~C.4]{ChenEtAl2023} that the slopes $(\Delta_\beta-\Delta_\alpha)/(\lambda_\beta-\lambda_\alpha)$ increase along a sequence of transitions.\label{ex:phases:hull}

\item \consolidation (The tempered posterior.) For the posterior at a fixed inverse temperature $\beta$ (\cref{ch:temperature}), show under \cref{ass:phases:region} that the local free energy $-\log\int_U\varphi e^{-\beta nL_n}$ is $\beta nL_n(w_U)+\lambda_U\log(n\beta)-(m_U-1)\log\log(n\beta)$ up to $O_p(1)$. Deduce that at fixed $n$ heating the posterior moves it toward simpler regions; for the double-well model at $n=10^4$ and $\delta=0.05$ the population switch is at $\beta=0.32$.\label{ex:phases:beta}

\item \byhand (Three regions.) Regions have $(\lambda,\Delta)=(\frac12,3\epsilon)$, $(\frac34,\epsilon)$ and $(1,0)$, with $\epsilon=10^{-3}$. Using the leading terms, find the order in which they are preferred and the switch points (answers: $n=842$ and $n=1885$). Replace the middle region by $(\frac34,2\epsilon)$, and show that it is never preferred and that the posterior jumps from the first region to the last at $n=1179$.\label{ex:phases:three}

\item \byhand (The emblem.) Show from its Newton polygon (\cref{thm:newton:plane}) that the chapter's emblem $K=a^6+a^2b^2+b^6$ has $\lambda=\frac12$ and $m=2$, and that region $A$ of the double-well model at $\delta=0$, where $K=\frac12a^2b^2+\frac12b^4+O(b^6)$, has the same values (compare $K$ with $a^2b^2+b^4$). Show that the constant $-\log(\sqrt{\pi/2}\,\phasesasinh(1/\sqrt{2\delta}))$ of $\bar F_A$ behaves like $-\log\log\frac1\delta$ as $\delta\to0$.\label{ex:phases:emblem}

\item \byhand\hard (A circle of equally good distributions.) Let $y=(\cos\theta)x_1+(\sin\theta)x_2+\varepsilon$ with $\theta$ uniform on $[0,2\pi)$, the truth $y=\varepsilon$, and a design with $\frac1n\sum_ix_ix_i^\top=I_2$ exactly. Show that every $\theta$ has misfit $\frac12$, that relatively finite variance fails, and, with $I_0$ the modified Bessel function and $Z_n=n^{-1/2}\sum_ix_i\varepsilon_i$, that
\[
F_n-nS_n=\tfrac n2-\log I_0\bigl(\sqrt n\,\abs{Z_n}\bigr)=\tfrac n2-\sqrt n\,\abs{Z_n}+\tfrac14\log n+O_p(1).
\]
Which conclusions of \cref{thm:phases:local} fail? \hint{$nL_n(\theta)-nS_n=\frac n2-\sqrt n\,\abs{Z_n}\cos(\theta-\theta_n)$ for an angle $\theta_n$.}\label{ex:phases:circle}

\item \variation (Why not the mixture.) For the mixture $(1-a)\Normal(0,1)+a\Normal(b,1)$ with a truth $(a_0,b_0)$, $0<a_0<1$, $b_0\ne0$, show that $\partial L/\partial a=-a_0(e^{bb_0}-1)$ at $(0,b)$, so that the arm $a=0$ of the cross is not a set of local minima where $bb_0>0$, and argue that no ridge separates the cross from the true point. What makes the segment of the double-well model different?\label{ex:phases:mixture}

\item \measurement (Sampling across the ridge.) Run the kit's random-walk Metropolis (\texttt{sk.metropolis}) on the posterior of the double-well model at $\delta=0.05$ and $n=10^4$: $64$ chains started on the segment $b=0$, steps tuned toward acceptance $0.3$ for $10^4$ discarded steps, then $10^5$ steps. How many chains reach $\abs b\ge\frac12$, although region $B$ has posterior mass near $0.97$? Relate the answer to $n\eta\approx140$. Then fix the step at $0.3$: why do the chains now cross, and at what cost? Try tempering.\label{ex:phases:mcmc}

\item \research (Scaling near singular truths.) For a family of truths $c_0(\delta)$ approaching a singular truth, find conditions under which the population free energy is $\lambda_0\log n+\Psi(n^\gamma\delta)+o(1)$, and relate them to the Newton polygon of $K$ at the singular truth. A test case: the product model with the truth $c_0=\delta$, where only the multiplicity changes.\label{ex:phases:scaling}

\item \research (Antecedents.) Gradient descent started near the segment of the double-well model stays there, while the posterior leaves at $n\approx3000$: a Bayesian transition without a dynamical one. Design a two-parameter model with a dynamical transition without a Bayesian antecedent, or argue that none exists under stated conditions.\label{ex:phases:antecedent}
\end{exercises}

\notesandrefs
\Textcite{ChenEtAl2023}, for whom jumps of the posterior between regions are universal in singular learning theory, build their definitions on \textcite[\S7.6]{Watanabe2009}, ``Singular learning process'', \textcite[\S9.4]{Watanabe2018} and a 2020 talk of Watanabe's; they point to the second and to Carroll's master's thesis \parencite{Carroll2021} for transitions in a hyperparameter of the prior or the truth instead of $n$, and state the local free energy formula, the switch equation, internal model selection and the Bayesian antecedent hypothesis. The position paper of \textcite{PepinLehalleurEtAl2025} uses the same comparison of local free energies, and \textcite{HooglandEtAl2025}, followed in \cref{ch:stages}, divide the training of transformers into stages at critical points of the estimated local learning coefficient \parencite{LauEtAl2025}. The blog sequence of \textcite{Carroll2023} is an informal introduction to phases and transitions. On the dynamical side, the segment of the double-well model is one of the lines of critical points that a smaller perceptron leaves in a larger one \parencite{FukumizuAmari2000}, and \cref{ch:sgd} compares gradient descent with the posterior.

\Cref{thm:phases:local}, a localization of the sandwich of \cref{thm:freeenergy:linear} to regions that need not contain the truth, the crossover curves of \cref{prop:phases:crossover}, and the double-well model with its statistic $\omega_n$ are the book's own; we have not found them in the literature, although \textcite{WatanabeAmari2003} studied the same scaling for layered models.

\paragraph{Translating the literature.} \Textcite{ChenEtAl2023} write $W_\alpha$ for our regions $U_\alpha$ and call them phases, $w^*_\alpha$ for our $w_\alpha$, and $F_n(W_\alpha)$ for the local free energy, with the prior multiplied by a partition of unity and renormalized, so that our $F_n(U_\alpha)$ is their $F_n(W_\alpha)+v_\alpha$; $\lambda_\alpha$ and $m_\alpha$ are ours. In their switch equation $n\Delta L+\Delta\lambda\log n+\Delta c=0$, $\Delta L$ is a difference of their losses $\E\norm{x-f}^2$, twice our $\Delta_\beta-\Delta_\alpha$ for the likelihood of their eq.~(3), and $c_\alpha$ collects $v_\alpha$ and the lower-order terms; $n_{\mathrm{cr}}$ is the critical sample size and $\alpha\to\beta$ a local phase transition. Their toy model writes $W$ and $b$ for our $V$ and $c$. Watanabe's 2009 book calls the free energy the stochastic complexity.

\printchapterbib
\end{refsection}
